% Standalone abstract for the public-capital-depreciation reframing.
% Used via % Standalone abstract for the public-capital-depreciation reframing.
% Used via % Standalone abstract for the public-capital-depreciation reframing.
% Used via % Standalone abstract for the public-capital-depreciation reframing.
% Used via \input{abstract.tex} from Article.tex inside the abstract environment.

\noindent Do housing markets efficiently price the physical depreciation of public capital? We investigate this question using a quasi-experimental design centered on narrow-margin referendums to renew local road-maintenance property taxes in Ohio. Combining 3,184 close-margin elections (1991--2021) with a novel ConvNeXt-V2 computer-vision model fine-tuned to measure physical road condition from satellite imagery, we provide evidence on the capitalization of infrastructure disinvestment into property values. A failed renewal cuts road-maintenance funding by 11 percent at the moment of the vote. Physical road quality declines by 23 percent over the subsequent ten years, with statistical significance emerging in the four-year-ahead window. House prices, however, show no immediate response: the price gap between narrowly-failing and narrowly-passing jurisdictions opens only in the fourth year, the same year visible road deterioration becomes statistically detectable. The cumulative ten-year price decline is 9 percent (\$15{,}350 on average), against an average per-household tax saving whose perpetuity present value is approximately \$1{,}900 --- a capitalization gap of roughly eight to one. Under rational expectations, the full future amenity decline should have been capitalized at the moment of the vote; the data decisively reject this benchmark. We rationalize the delayed-capitalization pattern using a dynamic general-equilibrium model in which households Bayesian-update on a noisy signal of the public-capital stock, recovering the rational-expectations benchmark as a knife-edge restriction. The calibrated model matches both the dynamic price path and the cross-sectional heterogeneity by urban density and price quantile. The implied marginal value of public funds is $-1.4$, with the experience-based counterpart in $[-2.1, -1.6]$, identifying the road-maintenance tax cut as a welfare-destroying policy whose cost is concealed from voters by the visibility lag of physical decay.
\\
\vspace{0in}\\
\noindent\textbf{JEL Codes:} R42, H71, R10, G14\\ from Article.tex inside the abstract environment.

\noindent Do housing markets efficiently price the physical depreciation of public capital? We investigate this question using a quasi-experimental design centered on narrow-margin referendums to renew local road-maintenance property taxes in Ohio. Combining 3,184 close-margin elections (1991--2021) with a novel ConvNeXt-V2 computer-vision model fine-tuned to measure physical road condition from satellite imagery, we provide evidence on the capitalization of infrastructure disinvestment into property values. A failed renewal cuts road-maintenance funding by 11 percent at the moment of the vote. Physical road quality declines by 23 percent over the subsequent ten years, with statistical significance emerging in the four-year-ahead window. House prices, however, show no immediate response: the price gap between narrowly-failing and narrowly-passing jurisdictions opens only in the fourth year, the same year visible road deterioration becomes statistically detectable. The cumulative ten-year price decline is 9 percent (\$15{,}350 on average), against an average per-household tax saving whose perpetuity present value is approximately \$1{,}900 --- a capitalization gap of roughly eight to one. Under rational expectations, the full future amenity decline should have been capitalized at the moment of the vote; the data decisively reject this benchmark. We rationalize the delayed-capitalization pattern using a dynamic general-equilibrium model in which households Bayesian-update on a noisy signal of the public-capital stock, recovering the rational-expectations benchmark as a knife-edge restriction. The calibrated model matches both the dynamic price path and the cross-sectional heterogeneity by urban density and price quantile. The implied marginal value of public funds is $-1.4$, with the experience-based counterpart in $[-2.1, -1.6]$, identifying the road-maintenance tax cut as a welfare-destroying policy whose cost is concealed from voters by the visibility lag of physical decay.
\\
\vspace{0in}\\
\noindent\textbf{JEL Codes:} R42, H71, R10, G14\\ from Article.tex inside the abstract environment.

\noindent Do housing markets efficiently price the physical depreciation of public capital? We investigate this question using a quasi-experimental design centered on narrow-margin referendums to renew local road-maintenance property taxes in Ohio. Combining 3,184 close-margin elections (1991--2021) with a novel ConvNeXt-V2 computer-vision model fine-tuned to measure physical road condition from satellite imagery, we provide evidence on the capitalization of infrastructure disinvestment into property values. A failed renewal cuts road-maintenance funding by 11 percent at the moment of the vote. Physical road quality declines by 23 percent over the subsequent ten years, with statistical significance emerging in the four-year-ahead window. House prices, however, show no immediate response: the price gap between narrowly-failing and narrowly-passing jurisdictions opens only in the fourth year, the same year visible road deterioration becomes statistically detectable. The cumulative ten-year price decline is 9 percent (\$15{,}350 on average), against an average per-household tax saving whose perpetuity present value is approximately \$1{,}900 --- a capitalization gap of roughly eight to one. Under rational expectations, the full future amenity decline should have been capitalized at the moment of the vote; the data decisively reject this benchmark. We rationalize the delayed-capitalization pattern using a dynamic general-equilibrium model in which households Bayesian-update on a noisy signal of the public-capital stock, recovering the rational-expectations benchmark as a knife-edge restriction. The calibrated model matches both the dynamic price path and the cross-sectional heterogeneity by urban density and price quantile. The implied marginal value of public funds is $-1.4$, with the experience-based counterpart in $[-2.1, -1.6]$, identifying the road-maintenance tax cut as a welfare-destroying policy whose cost is concealed from voters by the visibility lag of physical decay.
\\
\vspace{0in}\\
\noindent\textbf{JEL Codes:} R42, H71, R10, G14\\ from Article.tex inside the abstract environment.

\noindent Do housing markets efficiently price the physical depreciation of public capital? We investigate this question using a quasi-experimental design centered on narrow-margin referendums to renew local road-maintenance property taxes in Ohio. Combining 3,184 close-margin elections (1991--2021) with a novel ConvNeXt-V2 computer-vision model fine-tuned to measure physical road condition from satellite imagery, we provide evidence on the capitalization of infrastructure disinvestment into property values. A failed renewal cuts road-maintenance funding by 11 percent at the moment of the vote. Physical road quality declines by 23 percent over the subsequent ten years, with statistical significance emerging in the four-year-ahead window. House prices, however, show no immediate response: the price gap between narrowly-failing and narrowly-passing jurisdictions opens only in the fourth year, the same year visible road deterioration becomes statistically detectable. The cumulative ten-year price decline is 9 percent (\$15{,}350 on average), against an average per-household tax saving whose perpetuity present value is approximately \$1{,}900 --- a capitalization gap of roughly eight to one. Under rational expectations, the full future amenity decline should have been capitalized at the moment of the vote; the data decisively reject this benchmark. We rationalize the delayed-capitalization pattern using a dynamic general-equilibrium model in which households Bayesian-update on a noisy signal of the public-capital stock, recovering the rational-expectations benchmark as a knife-edge restriction. The calibrated model matches both the dynamic price path and the cross-sectional heterogeneity by urban density and price quantile. The implied marginal value of public funds is $-1.4$, with the experience-based counterpart in $[-2.1, -1.6]$, identifying the road-maintenance tax cut as a welfare-destroying policy whose cost is concealed from voters by the visibility lag of physical decay.
\\
\vspace{0in}\\
\noindent\textbf{JEL Codes:} R42, H71, R10, G14\\